\documentclass[11pt]{article}
\usepackage{../putnamprep}
\begin{document}
\sessiontitle{Test 005: Lesson and solution}{Technique: polynomial roots \quad$\cdot$\quad files under \texttt{techniques/12\_polynomial\_roots}}

\section*{The technique}

\paragraph{The fact.} A nonzero polynomial of degree $d$ has at most $d$ roots. So a polynomial with infinitely many roots is the zero polynomial.

\paragraph{The way it gets used.} If two polynomials agree at infinitely many points, they are the same polynomial. Apply the fact to their difference.

\paragraph{Tiny example.} Suppose $P$ is a polynomial with $P(n) = n^2$ for every positive integer $n$. Then $Q(x) = P(x) - x^2$ is zero at $1, 2, 3, \dots$ That is infinitely many roots, so $Q$ is zero and $P(x) = x^2$ for all real $x$.

\paragraph{Reach for it when} the problem says ``find all polynomials.'' The plan is almost always the same three steps:
\begin{enumerate}
  \item guess the answer from small values;
  \item show $P$ agrees with the guess at infinitely many points;
  \item conclude they are equal.
\end{enumerate}

\paragraph{The move in this problem.} The equation turns one known value into another. If you know $P(a)$, you get $P(a^2 + 1)$. Starting from $P(0) = 0$ that produces a chain of values, and induction carries it on forever.

\section*{Small cases}
$x = 0$: $P(1) = P(0)^2 + 1 = 1$. \quad $x = 1$: $P(2) = P(1)^2 + 1 = 2$. \quad $x = 2$: $P(5) = P(2)^2 + 1 = 5$. \quad $x = 5$: $P(26) = 26$.

The guess is $P(x) = x$.

\section*{Solution}
\begin{answer}
Only $P(x) = x$.
\end{answer}
\begin{proof}
\emph{It works.} If $P(x) = x$ then $P(0) = 0$ and $P(x^2 + 1) = x^2 + 1 = P(x)^2 + 1$.

\emph{Nothing else works.} Let $P$ satisfy the conditions. Define $a_0 = 0$ and $a_{k+1} = a_k^2 + 1$, so the sequence begins $0, 1, 2, 5, 26$.

\emph{Claim: $P(a_k) = a_k$ for every $k \ge 0$.} By induction on $k$. For $k = 0$, $P(0) = 0$ is given. If $P(a_k) = a_k$, then putting $x = a_k$ in the equation gives
\[ P(a_{k+1}) = P(a_k^2 + 1) = P(a_k)^2 + 1 = a_k^2 + 1 = a_{k+1}. \]

\emph{The $a_k$ are all different.} For every real $t$, $t^2 - t + 1 = \left(t - \tfrac12\right)^2 + \tfrac34 > 0$. So $a_{k+1} - a_k = a_k^2 - a_k + 1 > 0$, and the sequence is strictly increasing.

So $Q(x) = P(x) - x$ is zero at the infinitely many different points $a_0, a_1, a_2, \dots$ A nonzero polynomial has only finitely many roots. Hence $Q$ is the zero polynomial, and $P(x) = x$.
\end{proof}

\section*{How this is graded}
\begin{itemize}
  \item ``Find all'' needs both halves: $P(x) = x$ works, and nothing else does. Missing the first half costs a point or two.
  \item The $a_k$ must be shown to be different. If they repeated, there would be only finitely many roots and no conclusion. This is the step most often skipped.
  \item Guessing $P(x) = x$ from $P(1) = 1$, $P(2) = 2$, $P(5) = 5$ with no further argument earns about $2$.
\end{itemize}
\end{document}
